\section{Conclusiones}

\setlength{\parindent}{1.5cm}

\indent Llevar a cabo este proyecto nos permitió comprender de manera práctica el papel de los
conjuntos PRIMERO y SIGUIENTE en la construcción de analizadores sintácticos predictivos. A
diferencia del análisis léxico —que reconoce tokens mediante autómatas finitos— el análisis
sintáctico predictivo requiere razonar sobre las derivaciones posibles de la gramática, y los
conjuntos PRIMERO y SIGUIENTE son la herramienta que hace posible tomar esas decisiones sin
necesidad de retroceso.

\indent La implementación del algoritmo iterativo de punto fijo nos exigió comprender a fondo las
dependencias entre los conjuntos de distintos no terminales: un cambio en PRIMERO de un símbolo
puede propagarse hacia los conjuntos SIGUIENTE de otros. Manejar correctamente estas dependencias
mediante la repetición controlada de las reglas de inferencia fue el reto central del proyecto.

\indent Desde el punto de vista del diseño de software, la separación en capas —modelos,
persistencia y presentación— facilitó la división del trabajo entre sub equipos y permitió
desarrollar y probar cada módulo de forma independiente antes de la integración. La interfaz
gráfica, construida con Tkinter, añadió valor al proyecto al hacer accesible el cálculo de los
conjuntos a cualquier usuario sin necesidad de conocer el código.

\indent En conjunto, el proyecto consolidó nuestra comprensión de la segunda fase del diseño de
compiladores y reforzó buenas prácticas de colaboración, modularidad y documentación que
continuaremos aplicando en etapas posteriores del curso.
